\documentclass[spanish,letterpaper,oneside]{article}
\usepackage[utf8]{inputenc}
\usepackage[spanish]{babel}
\usepackage{geometry}
\geometry{letterpaper, margin=2.5cm}
\usepackage{xcolor}
\usepackage{enumitem}
\usepackage{tcolorbox}

\begin{document}
\begin{titlepage}
    \centering
    \vspace*{2cm}
    {\Huge \textbf{Universidad Autónoma de Coahuila} \par}
    \vspace{1cm}
    {\Large Facultad de Medicina \par}
    \vspace{2cm}
    {\LARGE \textbf{Examen de Ingreso: Biología Celular y Microbiología (Premedicina)} \par}
    \vspace{1cm}
    {\Large Guía Autocontenida y Banco de Reactivos EXANI-II \par}
    \vspace{3cm}
    {\large Elaborado por: Aspirante a Médico Cirujano \par}
    \vfill
    {\large \today \par}
\end{titlepage}

\tableofcontents
\newpage

\section{Módulo Teórico: Citología y Microbiología de Alto Rendimiento}
\subsection{Teoría Celular y Tipos Celulares}
\begin{itemize}
    \item \textbf{Célula Procariota (Bacterias, Arqueas):} Sin núcleo verdadero (nucleoide). Sin organelos membranosos. Cromosoma circular único. Ribosomas pequeños (70S). Tienen Pared Celular (peptidoglicano en bacterias).
    \item \textbf{Célula Eucariota (Animal, Vegetal, Fungi):} Núcleo definido por doble membrana (envoltura nuclear). Organelos membranosos complejos. Cromosomas lineales múltiples. Ribosomas (80S).
\end{itemize}

\subsection{Organelos y Membrana Plasmática}
La membrana es semipermeable, descrita por el \textbf{modelo de mosaico fluido} de Singer y Nicolson (fosfolípidos con movimiento lateral, colesterol estabilizador y proteínas integrales).
\begin{itemize}
    \item \textbf{Transporte Pasivo (a favor de gradiente, sin ATP):} Ósmosis (movimiento de $), difusión simple (gases lipofílicos) y difusión facilitada (requiere proteína transportadora, ej. para glucosa).
    \item \textbf{Transporte Activo (en contra de gradiente, requiere ATP):} Primario (ej. Bomba de ^+/K^+$ ATPasa, saca 3 Na+ y mete 2 K+) y Secundario (cotransporte).
    \item \textbf{Retículo Endoplásmico (RE):} Liso (REL) para síntesis lipídica y detoxificación. Rugoso (RER) tachonado de ribosomas, sintetiza proteínas de secreción.
    \item \textbf{Aparato de Golgi:} Centro de distribución, empaquetado y glicosilación post-traduccional de proteínas.
    \item \textbf{Lisosomas:} Estómago celular (fagocitosis y apoptosis), pH ácido.
\end{itemize}

\subsection{Ciclo Celular y División Celular}
\begin{tcolorbox}[colback=purple!5!white,colframe=purple!50!black,title=Fases de la Mitosis vs Meiosis]
\textbf{Mitosis:} División ecuacional de células somáticas. 1 diploide (2n) $\rightarrow$ 2 diploides idénticas (2n).\\
1. \textit{Profase:} Cromatina condensa en cromosomas. Se forma el huso acromático.\\
2. \textit{Metafase:} Cromosomas al ecuador (placa metafásica).\\
3. \textit{Anafase:} Cromátidas hermanas se separan (polos opuestos).\\
4. \textit{Telofase:} Célula se estrangula, reaparece membrana nuclear.\\
\textbf{Meiosis:} División reduccional (exclusiva de gónadas para generar espermatozoides/óvulos). 1 célula (2n) $\rightarrow$ 4 células haploides genéticamente distintas (n). Presenta recombinación genética (Crossing-over) en la Profase I.
\end{tcolorbox}

\subsection{Microbiología Básica}
\begin{itemize}
    \item \textbf{Virus:} Acelulares. Tienen ácido nucleico (ADN o ARN, nunca ambos) y cápside proteica. Son parásitos intracelulares estrictos.
    \item \textbf{Bacterias:} Procariotas. Se clasifican por su pared (Tinción Gram positiva = morada; Gram negativa = roja). Formas: Cocos, Bacilos, Espirilos.
    \item \textbf{Hongos:} Eucariotas saprófitos. Pared de quitina. Levaduras (unicelular) o Mohos (multicelular hifas).
    \item \textbf{Parásitos:} Eucariotas. Protozoarios (unicelulares móviles, ej. Amibas) y Helmintos (gusanos, metazoarios, ej. Tenia).
\end{itemize}

\newpage
\section{Banco de Reactivos: Simulación EXANI-II}
\begin{enumerate}[label=\textbf{\arabic*.}]
    \item \textbf{La teoría que sostiene que las mitocondrias y los cloroplastos se originaron a partir de bacterias procariotas que fueron fagocitadas por otra célula, estableciendo una relación simbiótica, se conoce como:}
    \begin{enumerate}
        \item[A)] Teoría Celular Clásica
        \item[B)] Teoría Endosimbiótica
        \item[C)] Generación Espontánea
        \item[D)] Panspermia
    \end{enumerate}
    \textit{Respuesta: B. Propuesta por Lynn Margulis, fundamentada en que mitocondrias/cloroplastos tienen ADN circular propio y ribosomas 70S.}

    \item \textbf{La bomba de Sodio-Potasio (Na+/K+ ATPasa) es un ejemplo biológico clásico de:}
    \begin{enumerate}
        \item[A)] Difusión simple
        \item[B)] Ósmosis
        \item[C)] Transporte pasivo facilitado
        \item[D)] Transporte activo primario
    \end{enumerate}
    \textit{Respuesta: D. Va en contra de los gradientes de concentración natural, requiriendo el uso de una molécula de ATP para funcionar.}

    \item \textbf{La enzima catalasa abundante en ciertas organelas descompone el peróxido de hidrógeno ($) en agua y oxígeno, protegiendo a la célula del estrés oxidativo. ¿Qué organelo es este?}
    \begin{enumerate}
        \item[A)] Peroxisoma
        \item[B)] Lisosoma
        \item[C)] Mitocondria
        \item[D)] Ribosoma
    \end{enumerate}
    \textit{Respuesta: A. Los peroxisomas son vesículas que contienen enzimas oxidativas (catalasa) dedicadas a la desintoxicación celular. Los lisosomas usan hidrolasas ácidas para digestión.}

    \item \textbf{¿Durante qué subfase específica del ciclo celular eucariota la célula detiene su división indefinidamente (latencia), siendo característico de las neuronas y miocitos adultos?}
    \begin{enumerate}
        \item[A)] Fase S
        \item[B)] Fase G0 (Quiescencia)
        \item[C)] Fase G2
        \item[D)] Citocinesis
    \end{enumerate}
    \textit{Respuesta: B. G0 es un estado donde la célula abandona el ciclo proliferativo, a veces permanentemente.}

    \item \textbf{Durante la meiosis se produce la variabilidad genética gracias al entrecruzamiento (crossing-over) de cromosomas homólogos. ¿En qué etapa exacta ocurre?}
    \begin{enumerate}
        \item[A)] Profase I
        \item[B)] Anafase I
        \item[C)] Profase II
        \item[D)] Metafase II
    \end{enumerate}
    \textit{Respuesta: A. El quiasma y el intercambio de alelos ocurre solo en la primera división meiótica (Paquiteno de Profase I).}

    \item \textbf{¿Qué componente de la pared celular es exclusivo de las bacterias y determina su tinción en la coloración de Gram?}
    \begin{enumerate}
        \item[A)] Quitina
        \item[B)] Celulosa
        \item[C)] Peptidoglicano (Mureína)
        \item[D)] Colesterol
    \end{enumerate}
    \textit{Respuesta: C. Una capa gruesa de peptidoglicano retiene el cristal violeta (Gram positivas), mientras que una fina con membrana externa se tiñe roja/rosada (Gram negativas).}

    \item \textbf{El eritrocito humano se coloca en una solución altamente salina (hipertónica). ¿Qué efecto celular observaremos a causa de la ósmosis?}
    \begin{enumerate}
        \item[A)] Plasmólisis o crenación (la célula se encoge)
        \item[B)] Lisis o turgencia (la célula explota)
        \item[C)] No habrá cambios de volumen
        \item[D)] Ingreso masivo de sodio
    \end{enumerate}
    \textit{Respuesta: A. El agua libre de la célula saldrá hacia el medio hipertónico intentando diluir la sal, causando que la célula se encoja.}

    \item \textbf{El Retículo Endoplásmico Liso (REL) a nivel celular, muy desarrollado en hepatocitos, es vital para:}
    \begin{enumerate}
        \item[A)] Traducción del ARN mensajero.
        \item[B)] Generar lisosomas primarios.
        \item[C)] Síntesis de lípidos (esteroides) y destoxificación hepática.
        \item[D)] Respiración celular aerobia.
    \end{enumerate}
    \textit{Respuesta: C. El REL carece de ribosomas y se encarga del metabolismo de carbohidratos, drogas/venenos y síntesis de hormonas lipídicas.}

    \item \textbf{Microorganismo unicelular eucariota, carente de pared celular, que se moviliza a través de pseudópodos, flagelos o cilios. Pertenece al reino Protista.}
    \begin{enumerate}
        \item[A)] Bacteria
        \item[B)] Hongo (levadura)
        \item[C)] Protozoario
        \item[D)] Virus
    \end{enumerate}
    \textit{Respuesta: C. Los protozoarios son eucariotas ameboideos o móviles. Bacterias carecen de núcleo. Hongos tienen pared celular.}

    \item \textbf{A nivel biológico, los virus no se consideran seres vivos de forma estricta porque:}
    \begin{enumerate}
        \item[A)] Tienen un núcleo difuso
        \item[B)] Contienen tanto ADN como ARN simultáneamente
        \item[C)] No tienen metabolismo propio ni maquinaria ribosómica
        \item[D)] Se replican por fisión binaria
    \end{enumerate}
    \textit{Respuesta: C. Los virus carecen de metabolismo (ribosomas, mitocondrias, producción de ATP) y deben forzar la maquinaria enzimática de la célula hospedera a reproducirlos.}

    \item \textbf{En un cariotipo normal de una célula somática humana (no gameto), ¿cuál es el número total de cromosomas?}
    \begin{enumerate}
        \item[A)] 23
        \item[B)] 46
        \item[C)] 22
        \item[D)] 92
    \end{enumerate}
    \textit{Respuesta: B. 46 cromosomas diploides (23 pares: 22 autosómicos y 1 par sexual). Los espermatozoides/óvulos tienen 23 (haploides).}

    \item \textbf{¿Qué proceso celular forma vesículas invaginando la membrana plasmática para incorporar líquidos o moléculas grandes hacia el interior de la célula?}
    \begin{enumerate}
        \item[A)] Exocitosis
        \item[B)] Endocitosis (Pinocitosis / Fagocitosis)
        \item[C)] Difusión simple
        \item[D)] Apoptosis
    \end{enumerate}
    \textit{Respuesta: B. La endocitosis permite meter partículas englobadas en membrana. Exocitosis es para secreción hacia afuera.}

    \item \textbf{El SARS-CoV-2 y el VIH comparten una característica biológica en común, y es que ambos son virus de tipo:}
    \begin{enumerate}
        \item[A)] Bacteriófagos
        \item[B)] Priones
        \item[C)] Virus cuyo genoma es de ARN
        \item[D)] Virus cuyo genoma es de ADN bicatenario
    \end{enumerate}
    \textit{Respuesta: C. El VIH es un retrovirus (ARN) y el coronavirus es de ARN de cadena positiva. Esto explica su altísima tasa de mutación a diferencia de los virus de ADN (ej. Herpes, VPH).}

    \item \textbf{Estructura celular responsable del anclaje y organización de los microtúbulos del huso acromático durante la división de células animales (los vegetales carecen de ella):}
    \begin{enumerate}
        \item[A)] Nucleolo
        \item[B)] Centrosoma (y centriolos)
        \item[C)] Complejo del Poro Nuclear
        \item[D)] Aparato de Golgi
    \end{enumerate}
    \textit{Respuesta: B. Los centriolos se replican y migran a los polos formando el centrosoma, origen del huso mitótico. (Las plantas organizan el huso sin centriolos verdaderos).}

    \item \textbf{A nivel estructural, la pared celular de los hongos verdaderos está compuesta bioquímicamente por:}
    \begin{enumerate}
        \item[A)] Celulosa
        \item[B)] Peptidoglicano
        \item[C)] Ácido hialurónico
        \item[D)] Quitina
    \end{enumerate}
    \textit{Respuesta: D. Quitina en los hongos. Celulosa en las plantas. Peptidoglicano en bacterias.}

    \item \textbf{Punto de control (checkpoint) más riguroso del ciclo celular eucariota (G1/S), en el cual el p53 revisa que el ADN no esté dañado antes de permitir su duplicación:}
    \begin{enumerate}
        \item[A)] Punto de restricción G1
        \item[B)] Punto M
        \item[C)] Control del Huso
        \item[D)] Punto de terminación telomérica
    \end{enumerate}
    \textit{Respuesta: A. El punto final de la fase G1 determina si la célula compromete recursos para dividirse. Si hay daño, la proteína p53 activa la apoptosis celular.}

    \item \textbf{En el modelo del Mosaico Fluido celular, ¿cuál molécula se inserta entre los fosfolípidos otorgando rigidez y disminuyendo la fluidez de la membrana a altas temperaturas?}
    \begin{enumerate}
        \item[A)] Glicoproteína
        \item[B)] Proteína periférica
        \item[C)] Colesterol
        \item[D)] Oligosacárido
    \end{enumerate}
    \textit{Respuesta: C. El colesterol es el amortiguador térmico de la membrana animal, modulando su consistencia mecánica.}

    \item \textbf{Tipo de asociación simbiótica en el que un microorganismo (ej. flora intestinal normal) se beneficia del huésped sin causarle daño ni otorgarle un beneficio directo importante:}
    \begin{enumerate}
        \item[A)] Mutualismo
        \item[B)] Parasitismo
        \item[C)] Comensalismo
        \item[D)] Depredación
    \end{enumerate}
    \textit{Respuesta: C. Comensalismo = +/0. (Mutualismo es +/+, Parasitismo es +/-).}

    \item \textbf{La citocinesis (separación del citoplasma) ocurre justo al finalizar la telofase de la mitosis. ¿En las células animales, esto se logra gracias a un anillo contráctil formado por?}
    \begin{enumerate}
        \item[A)] Celulosa
        \item[B)] Clatrina
        \item[C)] Actina y Miosina
        \item[D)] Microtúbulos (Tubulina)
    \end{enumerate}
    \textit{Respuesta: C. Los filamentos de actina y miosina crean un surco de segmentación. Las células vegetales (por tener pared) logran la separación formando un fragmoplasto (vesículas en el ecuador).}

    \item \textbf{Un microorganismo se clasifica como bacilo Gram-negativo. Esta denominación bacteriana nos indica morfológicamente que es:}
    \begin{enumerate}
        \item[A)] Forma esférica y de color morado
        \item[B)] Forma de bastón y de color rojo/rosa
        \item[C)] Forma espiral y de color azul
        \item[D)] Forma de coma y no se tiñe
    \end{enumerate}
    \textit{Respuesta: B. Bacilo = Bastón (cilíndrico). Gram negativo = se tiñe de fucsia/rosado/rojo con la safranina, pues perdió el cristal violeta de la pared delgada.}

\end{enumerate}
\end{document}
